
%% Concoct some colours of our own
\definecolor[named]{WTColorHeaderLineColor}{HTML}{FF0000}
\definecolor[named]{WTColorHeaderBackground}{HTML}{FFD5D2}
\definecolor[named]{WTColorFooterBackground}{HTML}{FF0000}
\definecolor[named]{WTColorChapterHeaderLightBackground}{HTML}{FFFFFF}

\usepackage{wallpaper}
\usepackage{changepage}
\usepackage[explicit]{titlesec}
\usepackage{fancyhdr}
\usepackage{tikz}
\usetikzlibrary{shapes,positioning}

\usepackage{geometry}
\geometry{
   paperwidth=216mm, paperheight=303mm,
   left=23mm,  %% or inner=23mm
   right=18mm, %% or outer=18mm
   top=23mm, bottom=23mm,
   headheight=\baselineskip,
   headsep=7mm,
   footskip=10mm
}


%% Command to hold chapter illustration image
\newcommand\chapterillustration{}

%% Define how the chapter title is printed
\titleformat{\chapter}{}{}{0pt}{
%% Background image at top of page
\ThisULCornerWallPaper{1}{resources/images/chapters/\chapterillustration}
%% Draw a semi-transparent rectangle across the top
\tikz[overlay,remember picture]
  \fill[WTColorChapterHeaderLightBackground,opacity=.7]
  (current page.north west) rectangle 
  ([yshift=-3cm] current page.north east);
  %% Check if on an odd or even page
  \strictpagecheck\checkoddpage
  %% On odd pages, "logo" image at lower right
  %% corner; Chapter number printed near spine
  %% edge (near the left); chapter title printed
  %% near outer edge (near the right).
  \ifoddpage{
    \ThisLRCornerWallPaper{.35}{resources/images/common/chapter_underlay}
    \begin{tikzpicture}[overlay,remember picture]
    \node[anchor=south west,
      xshift=20mm,yshift=-30mm,
      font=\sffamily\bfseries\huge] 
      at (current page.north west) 
      {\chaptername\ \thechapter};
    \node[fill=WTColorFooterBackground!80!black,text=white,
      font=\Huge\bfseries, 
      inner ysep=12pt, inner xsep=20pt,
      rounded rectangle,anchor=east, 
      xshift=-20mm,yshift=-30mm] 
      at (current page.north east) {#1}; % #1 ... chapter titel
    \end{tikzpicture}
  }
  %% On even pages, "logo" image at lower left
  %% corner; Chapter number printed near outer
  %% edge (near the right); chapter title printed
  %% near spine edge (near the left).
  \else {
    \ThisLLCornerWallPaper{.35}{resources/images/common/chapter_underlay}
    \begin{tikzpicture}[overlay,remember picture]
    \node[anchor=south east,
      xshift=-20mm,yshift=-30mm,
      font=\sffamily\bfseries\huge] 
      at (current page.north east)
      {\chaptername\ \thechapter};
    \node[fill=WTColorFooterBackground!80!black,text=white,
      font=\Huge\bfseries,
      inner sep=12pt, inner xsep=20pt,
      rounded rectangle,anchor=west,
      xshift=20mm,yshift=-30mm] 
      at ( current page.north west) {#1}; % #1 ... chapter titel
    \end{tikzpicture}
  }
  \fi
}
\titlespacing*{\chapter}{0pt}{0pt}{135mm}

%% Set the uniform width of the colour box
%% displaying the page number in footer
%% to the width of "99"
\newlength\pagenumwidth
\settowidth{\pagenumwidth}{99}

%% Define style of page number colour box
\tikzset{pagefooter/.style={
anchor=base,font=\sffamily\bfseries\small,
text=white,fill=WTColorFooterBackground!80!black,text centered,
text depth=17mm,text width=\pagenumwidth}}


%%%%%%%%%%
%%% Re-define running headers on non-chapter pages
%%%%%%%%%%



\fancypagestyle{headings}{%
  \fancyhf{}   % Clear all headers and footers first
  %% Right headers on odd pages
  \fancyhead[RO]{%
    %% First draw the background rectangles
    \begin{tikzpicture}[remember picture,overlay]
    \fill[WTColorHeaderBackground!25!white] (current page.north east) rectangle (current page.south west);
    \fill[white, rounded corners] ([xshift=-10mm,yshift=-20mm]current page.north east) rectangle ([xshift=15mm,yshift=17mm]current page.south west);
    \end{tikzpicture}
    %% Then the decorative line and the right mark
    \begin{tikzpicture}[xshift=-.75\baselineskip,yshift=.25\baselineskip,remember picture,    overlay,fill=WTColorHeaderLineColor,draw=WTColorHeaderLineColor]\fill circle(3pt);\draw[semithick](0,0) -- (current page.west |- 0,0);\end{tikzpicture} \sffamily\itshape\small\nouppercase{\rightmark}%\rightmark; \thesection=1.2  \sectionname
  }

  %% Left headers on even pages
  \fancyhead[LE]{%
    %% Background rectangles first
    \begin{tikzpicture}[remember picture,overlay]
    \fill[WTColorHeaderBackground!25!white] (current page.north east) rectangle (current page.south west);
    \fill[white, rounded corners] ([xshift=-15mm,yshift=-20mm]current page.north east) rectangle ([xshift=10mm,yshift=17mm]current page.south west);
    \end{tikzpicture}
    %% Then the right mark and the decorative line
    \sffamily\itshape\small\nouppercase{\leftmark}\ 
    \begin{tikzpicture}[xshift=.5\baselineskip,yshift=.25\baselineskip,remember picture, overlay,fill=WTColorHeaderLineColor,draw=WTColorHeaderLineColor]\fill (0,0) circle (3pt); \draw[semithick](0,0) -- (current page.east |- 0,0 );\end{tikzpicture}
  }

  %% Right footers on odd pages and left footers on even pages,
  %% display the page number in a colour box
  \fancyfoot[RO,LE]{\tikz[baseline]\node[pagefooter]{\thepage};}
  \renewcommand{\headrulewidth}{0pt}
  \renewcommand{\footrulewidth}{0pt}
}


%%%%%%%%%%
%%% Re-define running headers on chapter pages
%%%%%%%%%%
\fancypagestyle{plain}{%
  %% Clear all headers and footers
  \fancyhf{}
  %% Right footers on odd pages and left footers on even pages,
  %% display the page number in a colour box
  \fancyfoot[RO,LE]{\tikz[baseline]\node[pagefooter]{\thepage};}
  \renewcommand{\headrulewidth}{0pt}
  \renewcommand{\footrulewidth}{0pt}
}

%\renewcommand{\chaptermark}[1]{\markright{#1}{}}
%\renewcommand{\sectionmark}[1]{\markright{#1}{}}

%% Public domain image from
%% http://www.public-domain-image.com/objects/computer-chips/slides/six-computers-chips-circuits.html
%\renewcommand\chapterillustration{six-computers-chips-circuits}




